\documentclass[10pt,a4paper]{amsart}
\usepackage[margin=19mm]{geometry}
\usepackage{amsmath,amssymb,mathtools}
\usepackage[scr=rsfs,cal=euler]{mathalfa}
\usepackage{xfrac}
\usepackage[unicode,hidelinks,bookmarks=false]{hyperref}
\newcommand{\ie}{i.e.\ }
\newcommand{\cf}{cf.\ }
\newcommand{\resp}{respectively\ }
\newcommand{\Subsection}{\subsection}
\providecommand{\nobreakdash}{\nobreak\hbox{-}}
\setlength{\emergencystretch}{2em}
\multlinegap0pt
\usepackage{amssymb}
\usepackage{tikz}
\usepackage{enumerate}
\newtheorem{lemma}{Lemma}
\newcommand{\BNorm}[2]{\Big\|#1\Big\|_{#2}}
\newcommand{\Norm}[2]{\|#1\|_{#2}}
\title{$A_\infty$-invariance of oscillatory norms, and\\ Schatten characterisations of commutators}
\author{Tuomas Hyt\"onen}
\date{Source: arXiv:2604.22474v1 (CC BY 4.0)}
\begin{document}
\maketitle

\noindent\emph{Source and licence.} $A_\infty$-invariance of oscillatory norms, and\\ Schatten characterisations of commutators. Version arXiv:2604.22474v1 (CC BY 4.0), distributed by arXiv under the Creative Commons Attribution 4.0 International licence, http://creativecommons.org/licenses/by/4.0/, as recorded in the arXiv licence metadata for this exact version and shown on the abstract page https://arxiv.org/abs/2604.22474v1. Full licence notice: \texttt{licenses/arxiv-2604.22474-CC-BY-4.0.txt}.\par\smallskip

\noindent\textit{Editorial scope.} Counted unit: the article's lemma lem:Osc=Osc(D) (source0.tex lines 498--505). The article's complete original proof of it is delivered with it (source0.tex lines 507--535, one \texttt{proof} environment of 29 lines). No mathematics is written, repaired or re-derived: the statement and the proof are the article's own lines, in the article's own notation, and no retained line is substituted or re-flowed.  Retained uncounted as the source-local context the proof needs: lines 488--491.  Nothing was omitted from any retained span, so the package carries no omission record.  No retained line contains a citation, so the wrapper carries no bibliography and no bibliography entry.  No retained line contains a figure environment, an image-inclusion command or drawing code, so no image is delivered and the package declares no asset dependency.  Editorial formatting only, in addition: an AMS \texttt{amsart} preamble that restates the article's own declarations and macros used by the retained lines, namely the environments \texttt{lemma} on separate counters, together with the standard packages those lines need.  The retained environments print in this document as Lemma 1 in order of appearance, whereas the article prints the same items with the numbering of its own full document.  The complete native source of the article as distributed in the arXiv e-print for this exact version is bundled unchanged as \texttt{sources/199097/source0.tex}.
\begin{equation}\label{eq:adjacent}
  \forall\ B=B(x,t)\ \exists m\in\{1,\ldots,M\},\ \exists\ Q\in\mathscr D^m:\quad
  B\subset Q,\quad\operatorname{diam}(Q)\lesssim r.
\end{equation}

\begin{lemma}\label{lem:Osc=Osc(D)}
Let $(X,\rho,\mu)$ be a doubling metric measure space and $p,r\in(0,\infty)$ and $q\in(0,\infty]$.
For all measurable functions $f$ we have the following equivalence with implicit constants independent of $f$:
\begin{equation*}
  \Norm{f}{\operatorname{Osc}^{p,q}_r(\mu)}
  \sim\sum_{m=1}^M \Norm{f}{\operatorname{Osc}^{p,q}_r(\mathscr D^m,\mu)}
\end{equation*}
\end{lemma}

\begin{proof}
$\lesssim$: For every $Q\in \mathscr D$, we apply \eqref{eq:adjacent} to $B_Q$ to find a $Q^m\in\mathscr D^m$ with $B_Q\subset Q^m$ and $\operatorname{diam}(Q^m)\lesssim\ell(Q)$; hence $\mu(Q^m)\lesssim\mu(Q)\sim\mu(B_Q)$.
It follows that $\operatorname{osc}_{r,\mu}(f,B_Q)\lesssim\operatorname{osc}_{r,\mu}(f,Q^m)$. For $t'\neq t$, let $Q^{t'}:=\varnothing$ and $\operatorname{osc}_{r,\mu}(f,\varnothing):=0$.

Under the correspondence $Q\mapsto Q^m=R$, we claim that each $R\in\mathscr D^m$ may have at most boundedly many preimages. Indeed, these $Q$ satisfy $\ell(Q)\sim\ell(R)$, so they belong to boundedly many generations only. Within each fixed generation, these $Q$ are pairwise disjoint, contained in $R$, and their measure $\mu(Q)\gtrsim\mu(R)$ is at least a fixed fraction of $\mu(R)$. Hence, there can be only boundedly many such $Q$ in each generation, and then only boundedly many in the boundedly many generations altogether. It thus follows that
\begin{equation*}
\begin{split}
  \Norm{f}{\operatorname{Osc}^{p,q}_r(\mu)}
  &=\BNorm{\Big\{\operatorname{osc}_{r,\mu}(f;B_Q)\Big\}_{Q\in\mathscr D}}{\ell^{p,q}}
  \lesssim\sum_{m=1}^M \BNorm{\Big\{\operatorname{osc}_{r,\mu}(f;Q^m)\Big\}_{Q\in\mathscr D}}{\ell^{p,q}} \\
  &\lesssim\sum_{m=1}^M \BNorm{\Big\{\operatorname{osc}_{r,\mu}(f;R)\Big\}_{R\in\mathscr D^m}}{\ell^{p,q}}
  =\sum_{m=1}^M \Norm{f}{\operatorname{Osc}^{p,q}_r(\mathscr D^m,\mu)}.
\end{split}
\end{equation*}

$\gtrsim$: It suffices to consider a fixed $m$. For each $R\in\mathscr D^m$, there is $\hat R\in\mathscr D$ of the same side-length such that $z_R\in Q$. Then $R\subset B_{\hat R}$ when expansion factor defining the latter ball is fixed large enough, but also $\mu(B_{\hat R})\lesssim\mu(R)$. Hence $\operatorname{osc}_{r,\mu}(f;R)\lesssim\operatorname{osc}_{r,\mu}(f,B_{\hat R})$.

Reasoning as before, under the correspondence $R\mapsto \hat R=Q$, each $Q\in\mathscr D^m$ may have at most boundedly many preimages. Hence
\begin{equation*}
\begin{split}
  \Norm{f}{\operatorname{Osc}^{p,q}_r(\mathscr D^m,\mu)}
  &=\BNorm{\Big\{\operatorname{osc}_{r,\mu}(f;R)\Big\}_{R\in\mathscr D^m}}{\ell^{p,q}}
  \lesssim \BNorm{\Big\{\operatorname{osc}_{r,\mu}(f;B_{\hat R})\Big\}_{R\in\mathscr D^m}}{\ell^{p,q}} \\
  &\lesssim\BNorm{\Big\{\operatorname{osc}_{r,\mu}(f;B_Q)\Big\}_{Q\in\mathscr D}}{\ell^{p,q}}
  =\Norm{f}{\operatorname{Osc}^{p,q}_r(\mu)}.
\end{split}
\end{equation*}
This completes the proof.
\end{proof}

\end{document}
